\documentclass[
    language=english,
    doctype=handout,
    institution=none,
]{../../omnilatex}

\RequirePackage{config/document-settings}

\begin{document}
\maketitle

\section{Introduction}

Big-O notation describes the upper bound of an algorithm's time or space
complexity as a function of input size~$n$.  It answers the question:
\emph{how does runtime grow as input grows?}

Formally, $f(n) = O(g(n))$ if there exist constants $c > 0$ and
$n_0 \geq 0$ such that $f(n) \leq c \cdot g(n)$ for all $n \geq n_0$.

\keyconcept{Formal Definition}{%
    $f(n) = O(g(n))$ iff $\exists\, c > 0,\; n_0 \geq 0$ such that
    $0 \leq f(n) \leq c \cdot g(n)$ for all $n \geq n_0$.}

\keyconcept{Why Drop Constants and Lower Terms?}{%
    Big-O captures \emph{growth rate}, not exact counts.  Constants and
    lower-order terms become negligible as $n \to \infty$, so
    e.g.\ $3n^2 + 5n + 7 = O(n^2)$.}

\keyconcept{Best, Average, Worst Case}{%
    Big-O typically describes the \textbf{worst case}.  Omega~($\Omega$)
    gives the lower bound; Theta~($\Theta$) gives a tight bound
    (both upper and lower).}

\keyconcept{Amortized Analysis}{%
    Some operations are expensive occasionally but cheap on average.
    E.g.\ dynamic array append is $O(n)$ worst case but $O(1)$
    \emph{amortized} over a sequence of operations.}

\section{Common Complexities}

\begin{center}
\small
\begin{tabular}{@{}lcc@{}}
    \toprule
    \textbf{Complexity} & \textbf{Name} & \textbf{$n{=}1000$ ops} \\
    \midrule
    $O(1)$           & Constant     & 1           \\
    $O(\log n)$      & Logarithmic  & $\approx 10$\\
    $O(n)$           & Linear       & 1\,000      \\
    $O(n \log n)$    & Log-linear   & $\approx 10\,000$ \\
    $O(n^2)$         & Quadratic    & 1\,000\,000 \\
    $O(2^n)$         & Exponential  & $\approx 10^{301}$ \\
    $O(n!)$          & Factorial    & Astronomical \\
    \bottomrule
\end{tabular}
\end{center}

\remember{When comparing algorithms, prefer lower Big-O classes.
An $O(n \log n)$ sort always beats $O(n^2)$ for large inputs.}

\section{Quick Reference Rules}

\begin{itemize}
    \item \textbf{Loops:} Nested loops multiply --- $k$ nested loops $\Rightarrow O(n^k)$.
    \item \textbf{Sequential code:} Add complexities, keep the dominant term.
    \item \textbf{Divide and conquer:} $T(n) = aT(n/b) + O(n^d)$ --- use the Master Theorem.
    \item \textbf{Logarithmic factors} arise from halving/doubling (binary search, balanced trees).
\end{itemize}

\section{Practice Problems}

Determine the Big-O complexity of each:

\begin{enumerate}
    \item \textbf{Sum of elements:}
          \texttt{for i in range(n): s += a[i]}
    \item \textbf{Nested loops:}
          \texttt{for i in range(n):}\\
          \texttt{\ \ for j in range(n): print(a[i][j])}
    \item \textbf{Binary search} on a sorted array of size~$n$.
    \item \textbf{Merge sort:}
          $T(n) = 2T(n/2) + O(n)$.
    \item \textbf{Triple nested:}
          \texttt{for i in range(n):}\\
          \texttt{\ \ for j in range(i):}\\
          \texttt{\ \ \ \ for k in range(j): count += 1}
\end{enumerate}

\section{Further Reading}

\begin{itemize}
    \item Cormen, Leiserson, Rivest, Stein.
          \textit{Introduction to Algorithms}, 4th~ed., MIT Press, 2022.
    \item Sedgewick, Wayne.
          \textit{Algorithms}, 4th~ed., Addison-Wesley, 2011.
    \item Sipser, M.
          \textit{Introduction to the Theory of Computation}, 3rd~ed.,
          Cengage, 2012.
\end{itemize}

\end{document}
